\documentclass[instructions.tex]{subfiles}
\begin{document}
 
\chapter{De la fin de l'homme}
 
\primo Pourquoi suis-je au monde ? Ce n'est pas pour être riche et pour être honoré. Les honneurs et les biens de la terre ne sont pas dignes de moi, et ne peuvent me rendre content; je ne suis pas pour jouir des plaisirs des sens. Si cela était, l'homme dans sa fin n'aurait rien au dessus des animaux. Une créature raisonnable est créée pour une fin plus noble et plus digne de Dieu.

Je suis en ce monde pour y glorifier mon créateur, et pour me mettre en état de le glorifier et de le posséder en l'autre vie: voilà la fin que Dieu m'a destinée, et pour laquelle il m'a donné l'être; si je travaille pour une autre fin, je ne fais rien, et je ne mérite pas de vivre. Si le feu ne donnait aucune chaleur, il serait comme s'il n'était pas, parce qu'il serait inutile à la fin pour laquelle Dieu l'a créé. De même si l'homme ne s'appliquait pas à servir Dieu et à le glorifier, il serait sur la terre comme s'il n'était point: dès qu'il ne remplit pas la fin pour laquelle Dieu l'a mis au monde, il vaudrait mieux pour lui n'avoir jamais été; il devrait rentrer dans le néant, aussitôt qu'il ne tend pas à la fin pour laquelle Dieu l'a formé. 

Toutes les créatures visibles sont pour le service de l'homme, voilà leur fin: l'homme doit donc être et agir pour Dieu. Le créateur ne pouvait lui destiner une fin plus noble, et il ne pouvait lui en destiner une moindre. L'homme est très-honoré d'être destiné à une fin si glorieuse; mais il est bien malheureux de ne pas tendre à cette fin dans ses opérations. Un arbre planté dans une terre fertile pour donner des fruits à son maître, ne mérite que le feu quand il ne produit ni fruits ni feuilles. Quel châtiment ne mérite donc pas une créature, qui étant placée sur la terre par la main du Créateur pour l'honorer, ne fait rien pour sa gloire! 

\secundo Notre premier devoir est donc de glorifier Dieu sur la terre, en attendant que nous le glorifiions dans le ciel. Si nous lui demandons nos besoins et les choses temporelles, nous ne devons les demander qu'en vue de sa gloire. \og Vous n'êtes pas ici pour les biens de la terre\fg, disait sainte Térèse à ses religieuses, \og ni pour demander le succès des choses de ce monde, mais uniquement pour glorifier Dieu\fg. 

Cherchez donc, avant toutes choses, dit Jésus-Christ, le règne de Dieu\footnote{Quaerite ergò primùm regnum Dei. Matth.6.}. Un père doit faire régner Dieu dans sa famille; un pasteur, dans sa paroisse; un seigneur dans ses terres; un roi, dans ses états; c'est pour cette fin qu'ils ont l'autorité. Les rois vraiment grands et vraiment heureux, dit saint Augustin, sont ceux qui font usage de leur puissance pour soutenir le culte et la gloire de Dieu dans leurs empires\footnote{De Civ Deim l.5m c.24.}. 

Ce n'est point l'homme, qui n'est que misère, mais Dieu seul qui doit être glorifié\footnote{Soli Deo honor, gloria.}. Vous n'êtes point à vous, dit saint Augustin; vous êtes à Dieu, qui vous a fait ce que vous êtes; il est donc juste que vous viviez uniquement pour sa gloire\footnote{Tetum te exigit, qui totum te tecit. Aug.}. 

Qu'il est consolant et glorieux de vivre pour son Créateur! Si l'on se fait un honneur de servire les grands du monde, quel honneur de servir Dieu, le plus grand et le meilleur de tous les maitres! La moindre action faite pour sa gloire vaut incomparablement plus que tous les exploits des conquérants, et que tout ce qu'il y a de grand dans l'univers. Oh! que l'homme est ennemi de lui-même, quand il oublie ce qu'il doit à son auteur, puisqu'en oubliant ce qu'il doit à Dieu, il oublie son propre salut et son honneur!

\section{Résolutions}

\primo Puisque Dieu m'a créé pour sa gloire, je me proposerai désormais cette noble fin dans tout le détail de ma conduite. 

\secundo En conséquence, j'y rapporterai toutes mes actions et mes souffrances, ne manquant pas de les offrir à Dieu, dès le matin tous les jours, et renoncant à tout ce qui pourrait lui déplaire dans mes pensées, mes affections et mes démarches: offrande et renoncement que je renouvellerai de temps en temps pendant la journée, surtout quand j'aurai eu le malheur de les contredire par quelque faute notable. 

O mon Dieu ! je vous ai bien mal servi jusqu'ici; pardonnez-moi le passé, je vous en supplie, et faites-moi la grâce de m'attacher à vous comme à mon unique fin.

\end{document}
